\section{Introduction}
\noindent
Biological aging is the progressive accumulation of cellular damage leading to degeneration and organismal death \citep{aunan_molecular_2016}. DNA methylation patterns at CpG sites across the genome correlate strongly with the aging process, an effect that has been quantified using statistical models called ``methylation clocks'' \citep{jones_span_2015}. The first generation of methylation clocks were trained to predict chronological age from methylation levels at selected CpGs from across the genome \citep{horvath_dna_2013, hannum_genome-wide_2013, zhang_improved_2019}. A second generation of clocks were trained to use methylation levels to predict mortality risk as proxied by a combination of biomarkers of frailty and physiological decline \citep{levine_epigenetic_2018, lu_dna_2019}. Finally, a third generation of clocks have been trained to predict the rate of aging based on cohorts with longitudinal data on biomarkers of frailty~\citep{belsky_dunedinpace_2022}.

Greater predicted DNA methylation age compared to an individual's chronological age, known as methylation age acceleration, has been associated with an increased risk of many age-related diseases, including coronary heart disease, white matter hyperintensities, Type 2 diabetes mellitus, Parkinson's disease, and Alzheimer's disease (AD) \citep{horvath_epigenetic_2016, lu_dna_2019, raina_cerebral_2017, hodgson_epigenetic_2017, horvath_increased_2015, levine_epigenetic_2015, levine_epigenetic_2018}. As such, methylation clocks show potential as predictive biomarkers of the aging process and age-related health outcomes, and may capture relevant biological signals associated with aging. The clocks are also increasingly being used in social epidemiology research to quantify associations of methylation aging with exposure to adverse social and environmental factors that often differ across groups~\citep{non_social_2021, aiello_familial_2024, chiu_essential_2024, krieger_epigenetic_2024}.

While methylation is shaped by the environment of an individual, it is also strongly influenced by genetic variation \citep{kader_dna_2017}. Millions of methylation quantitative trait loci (meQTLs)---genetic variants that associate with the methylation level of a CpG site across individuals---have been identified \citep{smith_methylation_2014}. MeQTL influence methylation levels via many mechanisms, including disruption of CpGs and effects on transcription factor binding, gene expression, and other gene regulatory processes \citep{oliva_dna_2023, banovich_methylation_2014}. Methylation patterns also vary between human groups, and approximately 75\% of variance in methylation between human groups associates with genetic ancestry \citep{galanter_differential_2017}. This suggests that methylation levels and meQTLs often vary in frequency in different genetic ancestries. 

Despite these factors that lead to differential methylation levels in different genetic ancestries, the sociodemographic characteristics of the participants whose data were used to construct multiple commonly used methylation clocks is not known \citep{watkins_epigenetic_2023}. Most methylation clock studies do not report population, ancestry, or even geographic descriptors. Given that available genomic data are strongly biased towards individuals of European ancestry \citep{sirugo2019missing, popejoy_genomics_2016}, we hypothesized that lack of genetic diversity in the training data of methylation clocks could limit their generalizability across global and admixed populations. Similar factors have posed challenges for the application of polygenic risk scores (PRS) across human groups; PRS models often rapidly decrease in accuracy when applied to individuals not represented in the training set \citep{martin_clinical_2019, prive_portability_2022, novembre_addressing_2022}.

To quantify whether current methylation clocks are generalizable across global populations, we analyzed data from MAGENTA, a diverse AD study which has generated blood methylation and genotyping data for 621 individuals from the Americas, including genetically admixed individuals from African American, Puerto Rican, Cuban, and Peruvian cohorts. We evaluated the accuracy of first-, second-, and third-generation methylation clocks at predicting age in these individuals, along with three replication cohorts of African ancestry and European ancestry individuals. We then evaluated whether age acceleration metrics from these clocks associate with AD risk, as they do in individuals of European ancestry. Then, to investigate factors that influence clock performance, we quantified methylation levels, genetic variant frequency, and meQTL patterns for clock CpGs across genetic ancestries. Our results highlight obstacles to the application of methylation clocks as biomarkers for precision medicine and epidemiology, but they also identify promising avenues for considering genetic diversity in the development, application, and interpretation of methylation clocks.

%compute the frequency of genetic variants that disrupt CpGs sites across populations. We also intersect clock CpG sites with differentially methylated CpGs in African ancestry individuals relative to European ancestry individuals. Finally, we compare the frequencies of meQTL that influence clock CpG sites from three sets of independent meQTLs across different genetic ancestries. 

\clearpage